%%%%%%%%%%%%%%%%%%%%%%%%%%%%%%%%%%%%%%%%%%%%%%%%%%%%%%%%%%%%
% 9.8 DIFFERENTIAL-ALGEBRAIC EQUATIONS
% The Composition of Advanced Mathematics
%%%%%%%%%%%%%%%%%%%%%%%%%%%%%%%%%%%%%%%%%%%%%%%%%%%%%%%%%%%%

\section{Differential-Algebraic Equations}
\label{sec:differential-algebraic-equations}

\prerequisite{
Ordinary differential equations, systems of ODEs, linear algebra,
matrices, rank, eigenvalues, implicit differentiation, numerical methods,
dynamical systems, and basic constrained modeling.
}

\learningobjectives{
By the end of this section, the reader should be able to:
\begin{itemize}
    \item define a differential-algebraic equation;
    \item distinguish DAEs from ordinary differential equations;
    \item identify differential and algebraic variables;
    \item understand algebraic constraints;
    \item write DAEs in implicit form;
    \item recognize semi-explicit DAEs;
    \item identify consistent initial conditions;
    \item understand why arbitrary initial data may violate constraints;
    \item differentiate algebraic constraints when necessary;
    \item understand the idea of differentiation index;
    \item distinguish index-1 and higher-index systems conceptually;
    \item understand hidden constraints;
    \item recognize singular mass-matrix formulations;
    \item analyze linear constant-coefficient DAEs;
    \item distinguish regular and singular matrix pencils conceptually;
    \item understand constrained mechanical systems;
    \item formulate Lagrange-multiplier DAEs;
    \item recognize electrical-circuit DAEs;
    \item recognize chemical and process-system DAEs;
    \item understand conservation constraints;
    \item understand reduction of index-1 DAEs to ODEs locally;
    \item understand why higher-index DAEs are numerically difficult;
    \item recognize constraint drift;
    \item understand projection methods conceptually;
    \item understand implicit Euler for DAEs;
    \item recognize backward differentiation formulas;
    \item understand Newton iteration in implicit DAE solvers;
    \item understand event and constraint handling conceptually;
    \item recognize differential-algebraic boundary and control systems;
    \item connect DAEs with constrained dynamics, numerical analysis,
          optimization, and engineering models.
\end{itemize}
}

%%%%%%%%%%%%%%%%%%%%%%%%%%%%%%%%%%%%%%%%%%%%%%%%%%%%%%%%%%%%
\subsection{What Is a Differential-Algebraic Equation?}
\label{subsec:what-is-dae}
%%%%%%%%%%%%%%%%%%%%%%%%%%%%%%%%%%%%%%%%%%%%%%%%%%%%%%%%%%%%

A differential-algebraic equation, abbreviated DAE, combines differential
equations with algebraic constraints.

A simple example is

\[
\boxed{
x'(t)=y(t),
}
\]

together with

\[
\boxed{
x(t)^2+y(t)^2=1.
}
\]

The first equation describes differential evolution.

The second equation restricts the state to a constraint set.

%%%%%%%%%%%%%%%%%%%%%%%%%%%%%%%%%%%%%%%%%%%%%%%%%%%%%%%%%%%%
\subsection{ODE Versus DAE}
\label{subsec:ode-versus-dae}
%%%%%%%%%%%%%%%%%%%%%%%%%%%%%%%%%%%%%%%%%%%%%%%%%%%%%%%%%%%%

An explicit ODE has form

\[
\boxed{
\mathbf{x}'
=
\mathbf{f}(t,\mathbf{x}).
}
\]

A DAE may instead have implicit form

\[
\boxed{
\mathbf{F}
(
t,
\mathbf{x},
\mathbf{x}'
)
=
\mathbf{0}.
}
\]

The derivative may not be uniquely isolated for every component.

Some variables may be determined primarily by algebraic constraints.

%%%%%%%%%%%%%%%%%%%%%%%%%%%%%%%%%%%%%%%%%%%%%%%%%%%%%%%%%%%%
\subsection{Implicit Differential Equations}
\label{subsec:implicit-differential-equations}
%%%%%%%%%%%%%%%%%%%%%%%%%%%%%%%%%%%%%%%%%%%%%%%%%%%%%%%%%%%%

An equation such as

\[
\boxed{
F(t,x,x')=0
}
\]

is not automatically a DAE.

If

\[
\frac{\partial F}{\partial x'}
\]

is nonsingular in the relevant region, the implicit function theorem may
allow us to solve locally for

\[
\boxed{
x'=f(t,x).
}
\]

The system then behaves locally like an ODE.

A genuine DAE typically involves a singular derivative structure.

%%%%%%%%%%%%%%%%%%%%%%%%%%%%%%%%%%%%%%%%%%%%%%%%%%%%%%%%%%%%
\subsection{Semi-Explicit DAE Form}
\label{subsec:semi-explicit-dae}
%%%%%%%%%%%%%%%%%%%%%%%%%%%%%%%%%%%%%%%%%%%%%%%%%%%%%%%%%%%%

A common semi-explicit DAE has form

\[
\boxed{
\begin{aligned}
\mathbf{y}'
&=
\mathbf{f}
(
t,\mathbf{y},\mathbf{z}
),
\\
\mathbf{0}
&=
\mathbf{g}
(
t,\mathbf{y},\mathbf{z}
).
\end{aligned}
}
\]

Here:

\[
\mathbf{y}
=
\text{differential variables},
\]

and

\[
\mathbf{z}
=
\text{algebraic variables}.
\]

%%%%%%%%%%%%%%%%%%%%%%%%%%%%%%%%%%%%%%%%%%%%%%%%%%%%%%%%%%%%
\subsection{Differential Variables}
\label{subsec:differential-variables-dae}
%%%%%%%%%%%%%%%%%%%%%%%%%%%%%%%%%%%%%%%%%%%%%%%%%%%%%%%%%%%%

Differential variables have explicit time derivatives in the model.

For example,

\[
\boxed{
y'
=
f(t,y,z).
}
\]

Their values evolve dynamically.

%%%%%%%%%%%%%%%%%%%%%%%%%%%%%%%%%%%%%%%%%%%%%%%%%%%%%%%%%%%%
\subsection{Algebraic Variables}
\label{subsec:algebraic-variables-dae}
%%%%%%%%%%%%%%%%%%%%%%%%%%%%%%%%%%%%%%%%%%%%%%%%%%%%%%%%%%%%

Algebraic variables are constrained by equations such as

\[
\boxed{
g(t,y,z)=0.
}
\]

They may adjust instantaneously so the constraint remains satisfied.

This does not mean they are physically unimportant.

They may represent:

\[
\boxed{
\text{forces},
}
\]

\[
\boxed{
\text{pressures},
}
\]

\[
\boxed{
\text{currents},
}
\]

or

\[
\boxed{
\text{constraint multipliers}.
}
\]

%%%%%%%%%%%%%%%%%%%%%%%%%%%%%%%%%%%%%%%%%%%%%%%%%%%%%%%%%%%%
\subsection{Algebraic Constraint}
\label{subsec:dae-algebraic-constraint}
%%%%%%%%%%%%%%%%%%%%%%%%%%%%%%%%%%%%%%%%%%%%%%%%%%%%%%%%%%%%

An algebraic constraint restricts the allowed states.

For example,

\[
\boxed{
x^2+y^2=1
}
\]

restricts the state to the unit circle.

The DAE trajectory must remain on this constraint manifold.

%%%%%%%%%%%%%%%%%%%%%%%%%%%%%%%%%%%%%%%%%%%%%%%%%%%%%%%%%%%%
\subsection{Constraint Manifold}
\label{subsec:constraint-manifold-dae}
%%%%%%%%%%%%%%%%%%%%%%%%%%%%%%%%%%%%%%%%%%%%%%%%%%%%%%%%%%%%

Suppose

\[
\boxed{
g(\mathbf{x})=0.
}
\]

The set

\[
\boxed{
\mathcal{M}
=
\{
\mathbf{x}:
g(\mathbf{x})=0
\}
}
\]

is the constraint set.

When appropriate regularity conditions hold, it forms a manifold.

The DAE dynamics must remain tangent to \(\mathcal{M}\).

%%%%%%%%%%%%%%%%%%%%%%%%%%%%%%%%%%%%%%%%%%%%%%%%%%%%%%%%%%%%
\subsection{Consistent Initial Conditions}
\label{subsec:consistent-initial-conditions-dae}
%%%%%%%%%%%%%%%%%%%%%%%%%%%%%%%%%%%%%%%%%%%%%%%%%%%%%%%%%%%%

Initial values cannot usually be chosen arbitrarily.

They must satisfy all algebraic constraints.

For

\[
g(y,z)=0,
\]

consistent initial data satisfy

\[
\boxed{
g(y_0,z_0)=0.
}
\]

Additional hidden derivative constraints may also need to be satisfied.

%%%%%%%%%%%%%%%%%%%%%%%%%%%%%%%%%%%%%%%%%%%%%%%%%%%%%%%%%%%%
\subsection{Worked Example: Consistent Initial Data}
\label{subsec:worked-consistent-initial-data}
%%%%%%%%%%%%%%%%%%%%%%%%%%%%%%%%%%%%%%%%%%%%%%%%%%%%%%%%%%%%

\begin{workedexamplebox}{Checking a Constraint}

Consider

\[
x'=y,
\]

with constraint

\[
x^2+y^2=1.
\]

Suppose

\[
x(0)=\frac35.
\]

Consistency requires

\[
\left(
\frac35
\right)^2
+
y(0)^2
=
1.
\]

Thus,

\[
y(0)^2
=
1-\frac9{25}
=
\frac{16}{25}.
\]

Therefore,

\begin{answerbox}
\[
\boxed{
y(0)=\pm\frac45.
}
\]
\end{answerbox}

Any other value violates the algebraic constraint.

\end{workedexamplebox}

%%%%%%%%%%%%%%%%%%%%%%%%%%%%%%%%%%%%%%%%%%%%%%%%%%%%%%%%%%%%
\subsection{Hidden Constraints}
\label{subsec:hidden-constraints-dae}
%%%%%%%%%%%%%%%%%%%%%%%%%%%%%%%%%%%%%%%%%%%%%%%%%%%%%%%%%%%%

An algebraic equation may imply additional conditions after
differentiation.

Suppose

\[
\boxed{
g(y,z)=0.
}
\]

Differentiating with respect to time gives

\[
\boxed{
g_y y'
+
g_z z'
=
0
}
\]

for an autonomous constraint.

These derivative relations are called hidden constraints when they are
not explicitly listed in the original system.

%%%%%%%%%%%%%%%%%%%%%%%%%%%%%%%%%%%%%%%%%%%%%%%%%%%%%%%%%%%%
\subsection{Differentiating a Constraint}
\label{subsec:differentiating-dae-constraint}
%%%%%%%%%%%%%%%%%%%%%%%%%%%%%%%%%%%%%%%%%%%%%%%%%%%%%%%%%%%%

Consider

\[
\boxed{
x^2+y^2=1.
}
\]

Differentiate:

\[
2xx'
+
2yy'
=
0.
\]

Thus,

\[
\boxed{
xx'+yy'=0.
}
\]

This states that the velocity is tangent to the circle.

%%%%%%%%%%%%%%%%%%%%%%%%%%%%%%%%%%%%%%%%%%%%%%%%%%%%%%%%%%%%
\subsection{Tangent-Space Interpretation}
\label{subsec:dae-tangent-space}
%%%%%%%%%%%%%%%%%%%%%%%%%%%%%%%%%%%%%%%%%%%%%%%%%%%%%%%%%%%%

If

\[
g(\mathbf{x})=0,
\]

then along a valid trajectory,

\[
\boxed{
\nabla g(\mathbf{x})^T
\mathbf{x}'
=
0.
}
\]

Therefore the velocity lies in the tangent space of the constraint
manifold.

The gradient

\[
\nabla g
\]

is normal to the constraint surface.

%%%%%%%%%%%%%%%%%%%%%%%%%%%%%%%%%%%%%%%%%%%%%%%%%%%%%%%%%%%%
\subsection{Index of a DAE}
\label{subsec:dae-index}
%%%%%%%%%%%%%%%%%%%%%%%%%%%%%%%%%%%%%%%%%%%%%%%%%%%%%%%%%%%%

The index of a DAE measures, roughly, how many times algebraic constraints
must be differentiated before an explicit ODE-like description can be
obtained.

Several formal index definitions exist.

Therefore the word \emph{index} should always be interpreted in the
context of the chosen DAE theory.

%%%%%%%%%%%%%%%%%%%%%%%%%%%%%%%%%%%%%%%%%%%%%%%%%%%%%%%%%%%%
\subsection{Differentiation Index}
\label{subsec:differentiation-index}
%%%%%%%%%%%%%%%%%%%%%%%%%%%%%%%%%%%%%%%%%%%%%%%%%%%%%%%%%%%%

The differentiation index is the minimum number of differentiations of
the algebraic equations required, under suitable assumptions, to obtain
an explicit differential system.

Conceptually:

\[
\boxed{
\text{fewer differentiations}
\rightarrow
\text{lower index}.
}
\]

Higher index usually means greater analytical and numerical difficulty.

%%%%%%%%%%%%%%%%%%%%%%%%%%%%%%%%%%%%%%%%%%%%%%%%%%%%%%%%%%%%
\subsection{Index-1 Semi-Explicit DAE}
\label{subsec:index-one-dae}
%%%%%%%%%%%%%%%%%%%%%%%%%%%%%%%%%%%%%%%%%%%%%%%%%%%%%%%%%%%%

Consider

\[
\begin{aligned}
y'
&=
f(t,y,z),
\\
0
&=
g(t,y,z).
\end{aligned}
\]

If

\[
\boxed{
g_z
}
\]

is nonsingular, the implicit function theorem allows

\[
\boxed{
z=h(t,y)
}
\]

locally.

Then

\[
\boxed{
y'
=
f
(
t,y,h(t,y)
).
}
\]

Such systems are standard examples of index-1 DAEs.

%%%%%%%%%%%%%%%%%%%%%%%%%%%%%%%%%%%%%%%%%%%%%%%%%%%%%%%%%%%%
\subsection{Worked Example: Index-1 Reduction}
\label{subsec:worked-index-one-reduction}
%%%%%%%%%%%%%%%%%%%%%%%%%%%%%%%%%%%%%%%%%%%%%%%%%%%%%%%%%%%%

\begin{workedexamplebox}{Eliminating an Algebraic Variable}

Consider

\[
y'
=
-y+z,
\]

with constraint

\[
0=y+z-1.
\]

The algebraic equation gives

\[
z=1-y.
\]

Substitute into the differential equation:

\[
y'
=
-y+(1-y).
\]

Thus,

\[
y'
=
1-2y.
\]

Therefore the DAE reduces to

\begin{answerbox}
\[
\boxed{
y'
=
1-2y,
\qquad
z=1-y.
}
\]
\end{answerbox}

\end{workedexamplebox}

%%%%%%%%%%%%%%%%%%%%%%%%%%%%%%%%%%%%%%%%%%%%%%%%%%%%%%%%%%%%
\subsection{Higher-Index DAEs}
\label{subsec:higher-index-daes}
%%%%%%%%%%%%%%%%%%%%%%%%%%%%%%%%%%%%%%%%%%%%%%%%%%%%%%%%%%%%

A higher-index DAE cannot be reduced by simply solving the algebraic
constraint for the algebraic variables.

Differentiating constraints may be necessary.

Higher-index systems arise naturally in:

\[
\boxed{
\text{constrained mechanics},
}
\]

\[
\boxed{
\text{electrical circuits},
}
\]

and

\[
\boxed{
\text{multibody systems}.
}
\]

%%%%%%%%%%%%%%%%%%%%%%%%%%%%%%%%%%%%%%%%%%%%%%%%%%%%%%%%%%%%
\subsection{Simple Higher-Index Structure}
\label{subsec:simple-higher-index-structure}
%%%%%%%%%%%%%%%%%%%%%%%%%%%%%%%%%%%%%%%%%%%%%%%%%%%%%%%%%%%%

Consider

\[
\boxed{
x'=v,
}
\]

\[
\boxed{
v'=\lambda,
}
\]

with algebraic constraint

\[
\boxed{
x=0.
}
\]

Differentiate the constraint:

\[
x'=0.
\]

Thus,

\[
v=0.
\]

Differentiate again:

\[
v'=0.
\]

Hence,

\[
\lambda=0.
\]

The algebraic condition produces hidden differential constraints.

%%%%%%%%%%%%%%%%%%%%%%%%%%%%%%%%%%%%%%%%%%%%%%%%%%%%%%%%%%%%
\subsection{Mass-Matrix Form}
\label{subsec:dae-mass-matrix-form}
%%%%%%%%%%%%%%%%%%%%%%%%%%%%%%%%%%%%%%%%%%%%%%%%%%%%%%%%%%%%

Many DAEs are written

\[
\boxed{
M(t,x)x'
=
f(t,x),
}
\]

where \(M\) is called the mass matrix.

If \(M\) is invertible, then

\[
\boxed{
x'
=
M^{-1}f,
}
\]

and the system is an ODE.

If \(M\) is singular, the system may be a DAE.

%%%%%%%%%%%%%%%%%%%%%%%%%%%%%%%%%%%%%%%%%%%%%%%%%%%%%%%%%%%%
\subsection{Constant Linear DAE}
\label{subsec:constant-linear-dae}
%%%%%%%%%%%%%%%%%%%%%%%%%%%%%%%%%%%%%%%%%%%%%%%%%%%%%%%%%%%%

A linear constant-coefficient DAE has form

\[
\boxed{
E x'
=
A x
+
f(t),
}
\]

where \(E\) may be singular.

If

\[
\det(E)\neq0,
\]

the system reduces to an ODE.

If

\[
\det(E)=0,
\]

algebraic constraints may be present.

%%%%%%%%%%%%%%%%%%%%%%%%%%%%%%%%%%%%%%%%%%%%%%%%%%%%%%%%%%%%
\subsection{Matrix Pencil}
\label{subsec:dae-matrix-pencil}
%%%%%%%%%%%%%%%%%%%%%%%%%%%%%%%%%%%%%%%%%%%%%%%%%%%%%%%%%%%%

For the homogeneous linear DAE

\[
E x'=Ax,
\]

the matrix pencil is

\[
\boxed{
\lambda E-A.
}
\]

It generalizes the characteristic matrix

\[
\lambda I-A
\]

from ordinary linear systems.

%%%%%%%%%%%%%%%%%%%%%%%%%%%%%%%%%%%%%%%%%%%%%%%%%%%%%%%%%%%%
\subsection{Regular Matrix Pencil}
\label{subsec:regular-matrix-pencil}
%%%%%%%%%%%%%%%%%%%%%%%%%%%%%%%%%%%%%%%%%%%%%%%%%%%%%%%%%%%%

A square matrix pencil

\[
\lambda E-A
\]

is called regular if

\[
\boxed{
\det(\lambda E-A)
}
\]

is not identically zero as a polynomial in \(\lambda\).

Otherwise it is singular.

Regular pencils provide a standard framework for well-structured linear
DAEs.

%%%%%%%%%%%%%%%%%%%%%%%%%%%%%%%%%%%%%%%%%%%%%%%%%%%%%%%%%%%%
\subsection{Finite and Infinite Eigenvalues}
\label{subsec:finite-infinite-eigenvalues-dae}
%%%%%%%%%%%%%%%%%%%%%%%%%%%%%%%%%%%%%%%%%%%%%%%%%%%%%%%%%%%%

When \(E\) is singular, the pencil can possess finite generalized
eigenvalues together with structure associated with eigenvalues at
infinity.

The finite eigenvalues govern differential modes.

The infinite-eigenvalue structure is related to algebraic constraints and
DAE index.

%%%%%%%%%%%%%%%%%%%%%%%%%%%%%%%%%%%%%%%%%%%%%%%%%%%%%%%%%%%%
\subsection{Generalized Eigenvalue Problem}
\label{subsec:dae-generalized-eigenvalue-problem}
%%%%%%%%%%%%%%%%%%%%%%%%%%%%%%%%%%%%%%%%%%%%%%%%%%%%%%%%%%%%

Trying

\[
x(t)=ve^{\lambda t}
\]

in

\[
E x'=Ax
\]

gives

\[
\lambda Ev=Av.
\]

Thus,

\[
\boxed{
Av
=
\lambda Ev.
}
\]

This is a generalized eigenvalue problem.

%%%%%%%%%%%%%%%%%%%%%%%%%%%%%%%%%%%%%%%%%%%%%%%%%%%%%%%%%%%%
\subsection{Simple Linear DAE Example}
\label{subsec:simple-linear-dae-example}
%%%%%%%%%%%%%%%%%%%%%%%%%%%%%%%%%%%%%%%%%%%%%%%%%%%%%%%%%%%%

Consider

\[
\boxed{
\begin{pmatrix}
1&0\\
0&0
\end{pmatrix}
\begin{pmatrix}
x'\\
y'
\end{pmatrix}
=
\begin{pmatrix}
-x+y\\
x+y-1
\end{pmatrix}.
}
\]

The equations are

\[
\boxed{
x'=-x+y,
}
\]

and

\[
\boxed{
0=x+y-1.
}
\]

The second row is purely algebraic.

%%%%%%%%%%%%%%%%%%%%%%%%%%%%%%%%%%%%%%%%%%%%%%%%%%%%%%%%%%%%
\subsection{Reduction of the Linear DAE Example}
\label{subsec:reduction-linear-dae-example}
%%%%%%%%%%%%%%%%%%%%%%%%%%%%%%%%%%%%%%%%%%%%%%%%%%%%%%%%%%%%

From

\[
x+y-1=0,
\]

we obtain

\[
y=1-x.
\]

Then

\[
x'
=
-x+(1-x)
=
1-2x.
\]

Thus,

\[
\boxed{
x'=1-2x,
}
\]

and

\[
\boxed{
y=1-x.
}
\]

%%%%%%%%%%%%%%%%%%%%%%%%%%%%%%%%%%%%%%%%%%%%%%%%%%%%%%%%%%%%
\subsection{Solution of the Linear DAE Example}
\label{subsec:solution-linear-dae-example}
%%%%%%%%%%%%%%%%%%%%%%%%%%%%%%%%%%%%%%%%%%%%%%%%%%%%%%%%%%%%

The ODE

\[
x'=1-2x
\]

has solution

\[
\boxed{
x(t)
=
\frac12
+
Ce^{-2t}.
}
\]

Therefore,

\[
\boxed{
y(t)
=
\frac12
-
Ce^{-2t}.
}
\]

The initial conditions must satisfy

\[
\boxed{
x(0)+y(0)=1.
}
\]

%%%%%%%%%%%%%%%%%%%%%%%%%%%%%%%%%%%%%%%%%%%%%%%%%%%%%%%%%%%%
\subsection{Constrained Mechanical Systems}
\label{subsec:constrained-mechanical-systems-dae}
%%%%%%%%%%%%%%%%%%%%%%%%%%%%%%%%%%%%%%%%%%%%%%%%%%%%%%%%%%%%

Mechanical systems often involve geometric constraints.

Let

\[
q(t)
\]

denote generalized coordinates.

A holonomic constraint has form

\[
\boxed{
g(q,t)=0.
}
\]

The equations of motion must evolve while preserving this constraint.

%%%%%%%%%%%%%%%%%%%%%%%%%%%%%%%%%%%%%%%%%%%%%%%%%%%%%%%%%%%%
\subsection{Lagrange Multipliers in Constrained Mechanics}
\label{subsec:dae-lagrange-multipliers}
%%%%%%%%%%%%%%%%%%%%%%%%%%%%%%%%%%%%%%%%%%%%%%%%%%%%%%%%%%%%

A constrained mechanical system may be written

\[
\boxed{
M(q)q''
=
F(q,q',t)
+
G(q,t)^T\lambda,
}
\]

with

\[
\boxed{
g(q,t)=0.
}
\]

Here,

\[
G
=
\frac{\partial g}{\partial q},
\]

and

\[
\lambda
\]

contains Lagrange multipliers representing constraint forces.

%%%%%%%%%%%%%%%%%%%%%%%%%%%%%%%%%%%%%%%%%%%%%%%%%%%%%%%%%%%%
\subsection{Pendulum in Cartesian Coordinates}
\label{subsec:pendulum-cartesian-dae}
%%%%%%%%%%%%%%%%%%%%%%%%%%%%%%%%%%%%%%%%%%%%%%%%%%%%%%%%%%%%

A pendulum of fixed length \(L\) can be described using Cartesian
coordinates \(x,y\).

The geometric constraint is

\[
\boxed{
x^2+y^2=L^2.
}
\]

Using Newton's laws,

\[
\boxed{
m x''
=
-\lambda x,
}
\]

\[
\boxed{
m y''
=
-mg-\lambda y.
}
\]

Together with the length constraint, these equations form a DAE.

%%%%%%%%%%%%%%%%%%%%%%%%%%%%%%%%%%%%%%%%%%%%%%%%%%%%%%%%%%%%
\subsection{Velocity Constraint for the Pendulum}
\label{subsec:pendulum-velocity-constraint}
%%%%%%%%%%%%%%%%%%%%%%%%%%%%%%%%%%%%%%%%%%%%%%%%%%%%%%%%%%%%

Differentiate

\[
x^2+y^2=L^2.
\]

Then

\[
2xx'
+
2yy'
=
0.
\]

Thus,

\[
\boxed{
xx'+yy'=0.
}
\]

The velocity must be tangent to the circle.

%%%%%%%%%%%%%%%%%%%%%%%%%%%%%%%%%%%%%%%%%%%%%%%%%%%%%%%%%%%%
\subsection{Acceleration Constraint for the Pendulum}
\label{subsec:pendulum-acceleration-constraint}
%%%%%%%%%%%%%%%%%%%%%%%%%%%%%%%%%%%%%%%%%%%%%%%%%%%%%%%%%%%%

Differentiate again:

\[
x'^2
+
xx''
+
y'^2
+
yy''
=
0.
\]

Hence,

\[
\boxed{
xx''
+
yy''
=
-
(x'^2+y'^2).
}
\]

This relation can be combined with the equations of motion to solve for
the multiplier \(\lambda\).

%%%%%%%%%%%%%%%%%%%%%%%%%%%%%%%%%%%%%%%%%%%%%%%%%%%%%%%%%%%%
\subsection{Constraint Forces}
\label{subsec:constraint-forces-dae}
%%%%%%%%%%%%%%%%%%%%%%%%%%%%%%%%%%%%%%%%%%%%%%%%%%%%%%%%%%%%

The Lagrange multiplier does not represent an independently evolving
state.

It enforces the constraint.

For the pendulum, it is related to the tension needed to keep the mass on
the circular path.

Thus algebraic variables often have clear physical meaning.

%%%%%%%%%%%%%%%%%%%%%%%%%%%%%%%%%%%%%%%%%%%%%%%%%%%%%%%%%%%%
\subsection{Index of Constrained Mechanical DAEs}
\label{subsec:mechanical-dae-index}
%%%%%%%%%%%%%%%%%%%%%%%%%%%%%%%%%%%%%%%%%%%%%%%%%%%%%%%%%%%%

The standard position-level constrained mechanical formulation often has
a higher differentiation index.

The position constraint must be differentiated to obtain velocity and
acceleration constraints before the multiplier can be determined
explicitly.

This is one reason constrained mechanical DAEs can be numerically
challenging.

%%%%%%%%%%%%%%%%%%%%%%%%%%%%%%%%%%%%%%%%%%%%%%%%%%%%%%%%%%%%
\subsection{Electrical Circuit DAEs}
\label{subsec:electrical-circuit-daes}
%%%%%%%%%%%%%%%%%%%%%%%%%%%%%%%%%%%%%%%%%%%%%%%%%%%%%%%%%%%%

Electrical circuits can produce DAEs when Kirchhoff's laws impose
algebraic constraints.

Dynamic components include:

\[
\boxed{
\text{capacitors},
}
\]

and

\[
\boxed{
\text{inductors}.
}
\]

Algebraic components include ideal resistive or network constraints.

%%%%%%%%%%%%%%%%%%%%%%%%%%%%%%%%%%%%%%%%%%%%%%%%%%%%%%%%%%%%
\subsection{Capacitor and Inductor Relations}
\label{subsec:capacitor-inductor-dae}
%%%%%%%%%%%%%%%%%%%%%%%%%%%%%%%%%%%%%%%%%%%%%%%%%%%%%%%%%%%%

For a capacitor,

\[
\boxed{
i_C
=
C
\frac{dv_C}{dt}.
}
\]

For an inductor,

\[
\boxed{
v_L
=
L
\frac{di_L}{dt}.
}
\]

Kirchhoff's current and voltage laws impose additional algebraic
relations among currents and voltages.

%%%%%%%%%%%%%%%%%%%%%%%%%%%%%%%%%%%%%%%%%%%%%%%%%%%%%%%%%%%%
\subsection{Kirchhoff Current Law}
\label{subsec:kirchhoff-current-law-dae}
%%%%%%%%%%%%%%%%%%%%%%%%%%%%%%%%%%%%%%%%%%%%%%%%%%%%%%%%%%%%

At a circuit node,

\[
\boxed{
\sum_j i_j=0.
}
\]

This is algebraic in the instantaneous currents.

Combined with differential component laws, it produces a DAE system.

%%%%%%%%%%%%%%%%%%%%%%%%%%%%%%%%%%%%%%%%%%%%%%%%%%%%%%%%%%%%
\subsection{Kirchhoff Voltage Law}
\label{subsec:kirchhoff-voltage-law-dae}
%%%%%%%%%%%%%%%%%%%%%%%%%%%%%%%%%%%%%%%%%%%%%%%%%%%%%%%%%%%%

Around a closed loop,

\[
\boxed{
\sum_j v_j=0.
}
\]

Again, the equation is algebraic.

Circuit simulation software therefore frequently solves DAEs rather than
explicit ODEs.

%%%%%%%%%%%%%%%%%%%%%%%%%%%%%%%%%%%%%%%%%%%%%%%%%%%%%%%%%%%%
\subsection{Chemical Equilibrium Constraints}
\label{subsec:chemical-equilibrium-dae}
%%%%%%%%%%%%%%%%%%%%%%%%%%%%%%%%%%%%%%%%%%%%%%%%%%%%%%%%%%%%

Some chemical models combine slow dynamic reactions with rapidly
equilibrating processes.

A reaction may evolve according to differential equations while
equilibrium relations impose algebraic equations such as

\[
\boxed{
K_{eq}
=
\frac{
[C]^c[D]^d
}{
[A]^a[B]^b
}.
}
\]

This produces differential-algebraic models.

%%%%%%%%%%%%%%%%%%%%%%%%%%%%%%%%%%%%%%%%%%%%%%%%%%%%%%%%%%%%
\subsection{Process Systems}
\label{subsec:process-system-daes}
%%%%%%%%%%%%%%%%%%%%%%%%%%%%%%%%%%%%%%%%%%%%%%%%%%%%%%%%%%%%

Chemical engineering models may contain:

\[
\boxed{
\text{mass balances},
}
\]

\[
\boxed{
\text{energy balances},
}
\]

\[
\boxed{
\text{equations of state},
}
\]

and

\[
\boxed{
\text{phase-equilibrium constraints}.
}
\]

The conservation laws generate differential equations while constitutive
relations may be algebraic.

%%%%%%%%%%%%%%%%%%%%%%%%%%%%%%%%%%%%%%%%%%%%%%%%%%%%%%%%%%%%
\subsection{Conservation Constraints}
\label{subsec:conservation-constraints-dae}
%%%%%%%%%%%%%%%%%%%%%%%%%%%%%%%%%%%%%%%%%%%%%%%%%%%%%%%%%%%%

Suppose several state variables satisfy

\[
\boxed{
c_1x_1+\cdots+c_nx_n=C.
}
\]

This conservation law restricts the trajectory to a lower-dimensional
set.

One may eliminate a variable explicitly, or retain the constraint and
solve the resulting DAE.

%%%%%%%%%%%%%%%%%%%%%%%%%%%%%%%%%%%%%%%%%%%%%%%%%%%%%%%%%%%%
\subsection{Power-System DAEs}
\label{subsec:power-system-daes}
%%%%%%%%%%%%%%%%%%%%%%%%%%%%%%%%%%%%%%%%%%%%%%%%%%%%%%%%%%%%

Electrical power networks frequently combine dynamic generator models
with algebraic network-flow equations.

A schematic form is

\[
\boxed{
x'=f(x,z),
}
\]

\[
\boxed{
0=g(x,z).
}
\]

Here \(x\) may contain machine states while \(z\) contains network voltage
and phase variables.

%%%%%%%%%%%%%%%%%%%%%%%%%%%%%%%%%%%%%%%%%%%%%%%%%%%%%%%%%%%%
\subsection{DAEs in Control Theory}
\label{subsec:dae-control-theory}
%%%%%%%%%%%%%%%%%%%%%%%%%%%%%%%%%%%%%%%%%%%%%%%%%%%%%%%%%%%%

Descriptor systems are often written

\[
\boxed{
E x'
=
Ax+Bu,
}
\]

with output

\[
\boxed{
y=Cx+Du.
}
\]

If \(E\) is singular, the system is a descriptor or DAE control system.

Its analysis combines generalized eigenvalues with controllability and
observability concepts.

%%%%%%%%%%%%%%%%%%%%%%%%%%%%%%%%%%%%%%%%%%%%%%%%%%%%%%%%%%%%
\subsection{Consistent Initialization}
\label{subsec:dae-consistent-initialization}
%%%%%%%%%%%%%%%%%%%%%%%%%%%%%%%%%%%%%%%%%%%%%%%%%%%%%%%%%%%%

A DAE solver generally requires initial values satisfying both explicit
and hidden constraints.

For a semi-explicit index-1 DAE,

\[
\begin{aligned}
y'&=f(t,y,z),\\
0&=g(t,y,z),
\end{aligned}
\]

one may choose \(y_0\) and solve

\[
\boxed{
g(t_0,y_0,z_0)=0
}
\]

for \(z_0\).

%%%%%%%%%%%%%%%%%%%%%%%%%%%%%%%%%%%%%%%%%%%%%%%%%%%%%%%%%%%%
\subsection{Inconsistent Initial Conditions}
\label{subsec:dae-inconsistent-initial-conditions}
%%%%%%%%%%%%%%%%%%%%%%%%%%%%%%%%%%%%%%%%%%%%%%%%%%%%%%%%%%%%

Suppose the constraint is

\[
x+y=1.
\]

The initial choice

\[
x(0)=1,
\qquad
y(0)=1
\]

is inconsistent because

\[
1+1\neq1.
\]

A solver cannot produce a trajectory satisfying the original DAE from
such data without first modifying the initial state.

%%%%%%%%%%%%%%%%%%%%%%%%%%%%%%%%%%%%%%%%%%%%%%%%%%%%%%%%%%%%
\subsection{Initialization as a Nonlinear Solve}
\label{subsec:dae-initialization-nonlinear-solve}
%%%%%%%%%%%%%%%%%%%%%%%%%%%%%%%%%%%%%%%%%%%%%%%%%%%%%%%%%%%%

For nonlinear constraints,

\[
g(t_0,y_0,z_0)=0,
\]

consistent initialization may itself require a nonlinear solver.

Newton's method is commonly used.

Thus solving a DAE can begin with solving an algebraic system.

%%%%%%%%%%%%%%%%%%%%%%%%%%%%%%%%%%%%%%%%%%%%%%%%%%%%%%%%%%%%
\subsection{Index Reduction}
\label{subsec:dae-index-reduction}
%%%%%%%%%%%%%%%%%%%%%%%%%%%%%%%%%%%%%%%%%%%%%%%%%%%%%%%%%%%%

Higher-index systems are often transformed into lower-index systems before
numerical integration.

Methods may involve:

\[
\boxed{
\text{differentiating constraints},
}
\]

\[
\boxed{
\text{eliminating variables},
}
\]

or

\[
\boxed{
\text{introducing additional equations}.
}
\]

This process is called index reduction.

%%%%%%%%%%%%%%%%%%%%%%%%%%%%%%%%%%%%%%%%%%%%%%%%%%%%%%%%%%%%
\subsection{Danger of Naive Constraint Differentiation}
\label{subsec:danger-naive-index-reduction}
%%%%%%%%%%%%%%%%%%%%%%%%%%%%%%%%%%%%%%%%%%%%%%%%%%%%%%%%%%%%

Differentiating a constraint may enlarge the solution set.

For example,

\[
g(x)=0
\]

implies

\[
\frac{d}{dt}g(x(t))=0,
\]

but the derivative equation alone only guarantees

\[
g(x(t))
=
\text{constant},
\]

not necessarily zero.

Therefore the original constraint must still be enforced.

%%%%%%%%%%%%%%%%%%%%%%%%%%%%%%%%%%%%%%%%%%%%%%%%%%%%%%%%%%%%
\subsection{Constraint Drift}
\label{subsec:constraint-drift}
%%%%%%%%%%%%%%%%%%%%%%%%%%%%%%%%%%%%%%%%%%%%%%%%%%%%%%%%%%%%

A numerical method may approximately satisfy derivative constraints while
slowly moving away from the original algebraic constraint.

This is called constraint drift.

For example, a numerical pendulum may gradually violate

\[
\boxed{
x^2+y^2=L^2.
}
\]

%%%%%%%%%%%%%%%%%%%%%%%%%%%%%%%%%%%%%%%%%%%%%%%%%%%%%%%%%%%%
\subsection{Projection onto the Constraint Manifold}
\label{subsec:dae-projection}
%%%%%%%%%%%%%%%%%%%%%%%%%%%%%%%%%%%%%%%%%%%%%%%%%%%%%%%%%%%%

One approach to constraint drift is projection.

After a numerical step produces a tentative state

\[
\widetilde{x}_{n+1},
\]

replace it with a nearby corrected state \(x_{n+1}\) satisfying

\[
\boxed{
g(x_{n+1})=0.
}
\]

The correction projects the numerical solution back onto the constraint
manifold.

%%%%%%%%%%%%%%%%%%%%%%%%%%%%%%%%%%%%%%%%%%%%%%%%%%%%%%%%%%%%
\subsection{Baumgarte Stabilization Preview}
\label{subsec:baumgarte-stabilization-preview}
%%%%%%%%%%%%%%%%%%%%%%%%%%%%%%%%%%%%%%%%%%%%%%%%%%%%%%%%%%%%

For mechanical constraints

\[
g(q)=0,
\]

one may modify the differentiated constraint by adding feedback terms.

A schematic form is

\[
\boxed{
\ddot g
+
2\alpha\dot g
+
\beta^2g
=
0.
}
\]

This encourages constraint errors to decay.

This technique is called Baumgarte stabilization.

%%%%%%%%%%%%%%%%%%%%%%%%%%%%%%%%%%%%%%%%%%%%%%%%%%%%%%%%%%%%
\subsection{Implicit Numerical Methods}
\label{subsec:implicit-methods-daes}
%%%%%%%%%%%%%%%%%%%%%%%%%%%%%%%%%%%%%%%%%%%%%%%%%%%%%%%%%%%%

DAEs are usually solved with implicit methods.

At each time step, the new state must satisfy both:

\[
\boxed{
\text{the discretized differential equations},
}
\]

and

\[
\boxed{
\text{the algebraic constraints}.
}
\]

This produces an algebraic system at every step.

%%%%%%%%%%%%%%%%%%%%%%%%%%%%%%%%%%%%%%%%%%%%%%%%%%%%%%%%%%%%
\subsection{Implicit Euler for a Semi-Explicit DAE}
\label{subsec:implicit-euler-dae}
%%%%%%%%%%%%%%%%%%%%%%%%%%%%%%%%%%%%%%%%%%%%%%%%%%%%%%%%%%%%

Consider

\[
\begin{aligned}
y'&=f(t,y,z),\\
0&=g(t,y,z).
\end{aligned}
\]

Implicit Euler gives

\[
\boxed{
\frac{
y_{n+1}-y_n
}{
h
}
=
f
(
t_{n+1},
y_{n+1},
z_{n+1}
),
}
\]

together with

\[
\boxed{
0
=
g
(
t_{n+1},
y_{n+1},
z_{n+1}
).
}
\]

Both \(y_{n+1}\) and \(z_{n+1}\) are solved simultaneously.

%%%%%%%%%%%%%%%%%%%%%%%%%%%%%%%%%%%%%%%%%%%%%%%%%%%%%%%%%%%%
\subsection{Worked Example: Implicit Euler DAE Step}
\label{subsec:worked-implicit-euler-dae}
%%%%%%%%%%%%%%%%%%%%%%%%%%%%%%%%%%%%%%%%%%%%%%%%%%%%%%%%%%%%

\begin{workedexamplebox}{One Step of an Index-1 DAE}

Consider

\[
y'=-y+z,
\]

\[
0=y+z-1.
\]

Using step size \(h\),

\[
\frac{
y_{n+1}-y_n
}{
h
}
=
-y_{n+1}+z_{n+1}.
\]

The constraint gives

\[
z_{n+1}=1-y_{n+1}.
\]

Thus,

\[
\frac{
y_{n+1}-y_n
}{
h
}
=
1-2y_{n+1}.
\]

Multiply by \(h\):

\[
y_{n+1}-y_n
=
h-2hy_{n+1}.
\]

Therefore,

\[
(1+2h)y_{n+1}
=
y_n+h.
\]

Hence,

\begin{answerbox}
\[
\boxed{
y_{n+1}
=
\frac{
y_n+h
}{
1+2h
},
}
\]

and

\[
\boxed{
z_{n+1}
=
1-y_{n+1}.
}
\]
\end{answerbox}

\end{workedexamplebox}

%%%%%%%%%%%%%%%%%%%%%%%%%%%%%%%%%%%%%%%%%%%%%%%%%%%%%%%%%%%%
\subsection{Nonlinear Solve at Each Step}
\label{subsec:dae-nonlinear-solve}
%%%%%%%%%%%%%%%%%%%%%%%%%%%%%%%%%%%%%%%%%%%%%%%%%%%%%%%%%%%%

For nonlinear DAEs, implicit discretization produces equations

\[
\boxed{
R(w)=0,
}
\]

where \(w\) contains the unknown new-time variables.

Newton iteration may be used:

\[
\boxed{
w^{(k+1)}
=
w^{(k)}
-
J_R(w^{(k)})^{-1}
R(w^{(k)}).
}
\]

In practice, the linear Newton system is solved without explicitly forming
the inverse matrix.

%%%%%%%%%%%%%%%%%%%%%%%%%%%%%%%%%%%%%%%%%%%%%%%%%%%%%%%%%%%%
\subsection{Jacobian Structure in DAE Solvers}
\label{subsec:dae-jacobian-structure}
%%%%%%%%%%%%%%%%%%%%%%%%%%%%%%%%%%%%%%%%%%%%%%%%%%%%%%%%%%%%

For

\[
F(t,x,x')=0,
\]

an implicit time method often requires matrices involving

\[
\boxed{
F_x
}
\]

and

\[
\boxed{
F_{x'}.
}
\]

The combination may look like

\[
\boxed{
F_x
+
\alpha
F_{x'},
}
\]

where \(\alpha\) depends on the time discretization.

Efficient DAE solvers exploit sparsity in these matrices.

%%%%%%%%%%%%%%%%%%%%%%%%%%%%%%%%%%%%%%%%%%%%%%%%%%%%%%%%%%%%
\subsection{Backward Differentiation Formulas}
\label{subsec:bdf-daes}
%%%%%%%%%%%%%%%%%%%%%%%%%%%%%%%%%%%%%%%%%%%%%%%%%%%%%%%%%%%%

Backward differentiation formulas, abbreviated BDF methods, are widely
used for stiff ODEs and DAEs.

A first-order BDF is implicit Euler.

A second-order BDF approximates

\[
x'(t_{n+1})
\]

by

\[
\boxed{
\frac{
3x_{n+1}
-
4x_n
+
x_{n-1}
}{
2h
}.
}
\]

The approximation is inserted into

\[
F(t,x,x')=0.
\]

%%%%%%%%%%%%%%%%%%%%%%%%%%%%%%%%%%%%%%%%%%%%%%%%%%%%%%%%%%%%
\subsection{Why Explicit Euler Is Problematic for DAEs}
\label{subsec:explicit-euler-dae-problem}
%%%%%%%%%%%%%%%%%%%%%%%%%%%%%%%%%%%%%%%%%%%%%%%%%%%%%%%%%%%%

An explicit update may advance differential variables without ensuring
that the new state satisfies

\[
\boxed{
g(y_{n+1},z_{n+1})=0.
}
\]

The algebraic variables may not even have an explicit update rule.

Therefore implicit constraint-aware methods are generally preferred.

%%%%%%%%%%%%%%%%%%%%%%%%%%%%%%%%%%%%%%%%%%%%%%%%%%%%%%%%%%%%
\subsection{DAE Stiffness}
\label{subsec:dae-stiffness}
%%%%%%%%%%%%%%%%%%%%%%%%%%%%%%%%%%%%%%%%%%%%%%%%%%%%%%%%%%%%

DAEs frequently occur in models with widely separated time scales.

Examples include:

\[
\boxed{
\text{fast electrical transients},
}
\]

\[
\boxed{
\text{chemical equilibria},
}
\]

and

\[
\boxed{
\text{mechanical constraints}.
}
\]

Implicit methods are often chosen because they have favorable stability
properties for stiff systems.

%%%%%%%%%%%%%%%%%%%%%%%%%%%%%%%%%%%%%%%%%%%%%%%%%%%%%%%%%%%%
\subsection{Scaling of DAE Systems}
\label{subsec:dae-scaling}
%%%%%%%%%%%%%%%%%%%%%%%%%%%%%%%%%%%%%%%%%%%%%%%%%%%%%%%%%%%%

Poorly scaled equations can make nonlinear solves difficult.

Suppose one constraint has magnitude

\[
10^{-8}
\]

while another has magnitude

\[
10^6.
\]

A solver may treat residuals unevenly.

Nondimensionalization or explicit scaling can improve conditioning.

%%%%%%%%%%%%%%%%%%%%%%%%%%%%%%%%%%%%%%%%%%%%%%%%%%%%%%%%%%%%
\subsection{Residual Formulation}
\label{subsec:dae-residual-formulation}
%%%%%%%%%%%%%%%%%%%%%%%%%%%%%%%%%%%%%%%%%%%%%%%%%%%%%%%%%%%%

A general solver may work directly with

\[
\boxed{
F(t,x,x')=0.
}
\]

Rather than distinguishing every equation manually as differential or
algebraic, the solver seeks a state and derivative whose residual is near
zero.

This formulation is flexible for large coupled systems.

%%%%%%%%%%%%%%%%%%%%%%%%%%%%%%%%%%%%%%%%%%%%%%%%%%%%%%%%%%%%
\subsection{Error Control}
\label{subsec:dae-error-control}
%%%%%%%%%%%%%%%%%%%%%%%%%%%%%%%%%%%%%%%%%%%%%%%%%%%%%%%%%%%%

DAE error control must monitor both differential accuracy and algebraic
constraint satisfaction.

Useful quantities include

\[
\boxed{
\|F(t,x,x')\|,
}
\]

and

\[
\boxed{
\|g(t,y,z)\|.
}
\]

Small integration error does not automatically imply small constraint
error.

%%%%%%%%%%%%%%%%%%%%%%%%%%%%%%%%%%%%%%%%%%%%%%%%%%%%%%%%%%%%
\subsection{Events in DAE Models}
\label{subsec:dae-events}
%%%%%%%%%%%%%%%%%%%%%%%%%%%%%%%%%%%%%%%%%%%%%%%%%%%%%%%%%%%%

Real DAE models may switch structure when an event occurs.

Examples include:

\[
\boxed{
\text{contact},
}
\]

\[
\boxed{
\text{valve opening},
}
\]

\[
\boxed{
\text{electrical switching},
}
\]

or

\[
\boxed{
\text{phase change}.
}
\]

After an event, the algebraic constraints may change and consistent
initialization may need to be performed again.

%%%%%%%%%%%%%%%%%%%%%%%%%%%%%%%%%%%%%%%%%%%%%%%%%%%%%%%%%%%%
\subsection{Hybrid DAE Systems}
\label{subsec:hybrid-dae-systems}
%%%%%%%%%%%%%%%%%%%%%%%%%%%%%%%%%%%%%%%%%%%%%%%%%%%%%%%%%%%%

A hybrid DAE combines continuous constrained dynamics with discrete mode
changes.

A schematic representation is

\[
\boxed{
F_m(t,x,x')=0,
}
\]

where the active equation set depends on a discrete mode \(m\).

Switching can change:

\[
\boxed{
\text{constraints},
}
\]

\[
\boxed{
\text{parameters},
}
\]

or

\[
\boxed{
\text{system dimension}.
}
\]

%%%%%%%%%%%%%%%%%%%%%%%%%%%%%%%%%%%%%%%%%%%%%%%%%%%%%%%%%%%%
\subsection{DAEs and Optimization}
\label{subsec:dae-optimization-connection}
%%%%%%%%%%%%%%%%%%%%%%%%%%%%%%%%%%%%%%%%%%%%%%%%%%%%%%%%%%%%

Algebraic constraints are central to constrained optimization.

For example, the first-order conditions for minimizing

\[
f(x)
\]

subject to

\[
g(x)=0
\]

introduce Lagrange multipliers:

\[
\boxed{
\nabla f(x)
+
G(x)^T\lambda
=
0,
}
\]

\[
\boxed{
g(x)=0.
}
\]

Dynamic optimization combines these ideas with differential equations and
can produce differential-algebraic optimality systems.

%%%%%%%%%%%%%%%%%%%%%%%%%%%%%%%%%%%%%%%%%%%%%%%%%%%%%%%%%%%%
\subsection{DAEs and Constrained Dynamical Systems}
\label{subsec:dae-constrained-dynamical-systems}
%%%%%%%%%%%%%%%%%%%%%%%%%%%%%%%%%%%%%%%%%%%%%%%%%%%%%%%%%%%%

A DAE can be viewed as a dynamical system restricted to a constraint
manifold.

The dynamics must satisfy:

\[
\boxed{
\text{evolution equations},
}
\]

and

\[
\boxed{
\text{geometric or physical constraints}.
}
\]

This viewpoint connects DAEs with differential geometry and constrained
mechanics.

%%%%%%%%%%%%%%%%%%%%%%%%%%%%%%%%%%%%%%%%%%%%%%%%%%%%%%%%%%%%
\subsection{DAEs and Singular Perturbations}
\label{subsec:dae-singular-perturbation}
%%%%%%%%%%%%%%%%%%%%%%%%%%%%%%%%%%%%%%%%%%%%%%%%%%%%%%%%%%%%

Some DAEs arise as limits of ODE systems containing a small parameter.

For example,

\[
\boxed{
y'
=
f(y,z),
}
\]

\[
\boxed{
\varepsilon z'
=
g(y,z).
}
\]

As

\[
\varepsilon\to0,
\]

the second equation formally becomes

\[
\boxed{
0=g(y,z),
}
\]

producing a DAE.

%%%%%%%%%%%%%%%%%%%%%%%%%%%%%%%%%%%%%%%%%%%%%%%%%%%%%%%%%%%%
\subsection{Fast and Slow Variables}
\label{subsec:dae-fast-slow-variables}
%%%%%%%%%%%%%%%%%%%%%%%%%%%%%%%%%%%%%%%%%%%%%%%%%%%%%%%%%%%%

In

\[
\varepsilon z'
=
g(y,z),
\]

small \(\varepsilon\) means \(z\) can evolve on a much faster time scale
than \(y\).

The DAE limit assumes the fast variable has relaxed instantly to the
constraint

\[
\boxed{
g(y,z)=0.
}
\]

This is the basis of many quasi-steady-state approximations.

%%%%%%%%%%%%%%%%%%%%%%%%%%%%%%%%%%%%%%%%%%%%%%%%%%%%%%%%%%%%
\subsection{Differential Index Versus Numerical Difficulty}
\label{subsec:dae-index-numerical-difficulty}
%%%%%%%%%%%%%%%%%%%%%%%%%%%%%%%%%%%%%%%%%%%%%%%%%%%%%%%%%%%%

Higher index often correlates with numerical difficulty, but index alone
does not completely determine solver performance.

Other important factors include:

\[
\boxed{
\text{conditioning},
}
\]

\[
\boxed{
\text{stiffness},
}
\]

\[
\boxed{
\text{scaling},
}
\]

and

\[
\boxed{
\text{constraint geometry}.
}
\]

%%%%%%%%%%%%%%%%%%%%%%%%%%%%%%%%%%%%%%%%%%%%%%%%%%%%%%%%%%%%
\subsection{DAE Modeling Workflow}
\label{subsec:dae-modeling-workflow}
%%%%%%%%%%%%%%%%%%%%%%%%%%%%%%%%%%%%%%%%%%%%%%%%%%%%%%%%%%%%

A practical DAE workflow is:

\begin{enumerate}
    \item identify all dynamic variables;
    \item identify algebraic constraints;
    \item choose differential and algebraic variables;
    \item write the model in implicit or semi-explicit form;
    \item inspect the derivative matrix or mass matrix;
    \item determine whether the system reduces directly to an ODE;
    \item identify consistent initial conditions;
    \item determine whether hidden constraints exist;
    \item estimate or determine the DAE index when necessary;
    \item perform index reduction if required;
    \item select an implicit numerical solver;
    \item monitor algebraic residuals;
    \item monitor constraint drift;
    \item use projection or stabilization if appropriate;
    \item verify physical conservation laws;
    \item test sensitivity to time step and solver tolerances.
\end{enumerate}

%%%%%%%%%%%%%%%%%%%%%%%%%%%%%%%%%%%%%%%%%%%%%%%%%%%%%%%%%%%%
\subsection{Choosing a DAE Formulation}
\label{subsec:choosing-dae-formulation}
%%%%%%%%%%%%%%%%%%%%%%%%%%%%%%%%%%%%%%%%%%%%%%%%%%%%%%%%%%%%

If algebraic variables can be solved easily from

\[
g(t,y,z)=0,
\]

consider:

\[
\boxed{
\text{local reduction to an ODE}.
}
\]

If the constraint has strong physical meaning, consider retaining it as a
DAE.

For constrained mechanics, Lagrange multipliers often preserve physical
interpretation.

For large engineering systems, a residual form

\[
\boxed{
F(t,x,x')=0
}
\]

is often the most natural.

%%%%%%%%%%%%%%%%%%%%%%%%%%%%%%%%%%%%%%%%%%%%%%%%%%%%%%%%%%%%
\subsection{Common Errors}
\label{subsec:dae-common-errors}
%%%%%%%%%%%%%%%%%%%%%%%%%%%%%%%%%%%%%%%%%%%%%%%%%%%%%%%%%%%%

\subsubsection{Calling Every Implicit ODE a DAE}

An implicit equation

\[
F(t,x,x')=0
\]

may still be an ODE if it can be solved locally for \(x'\).

A DAE involves genuine algebraic or singular derivative structure.

%%%%%%%%%%%%%%%%%%%%%%%%%%%%%%%%%%%%%%%%%%%%%%%%%%%%%%%%%%%%
\subsubsection{Ignoring Algebraic Constraints at the Initial Time}

Initial data must satisfy

\[
\boxed{
g(t_0,y_0,z_0)=0.
}
\]

Otherwise the initial state does not lie on the admissible constraint
set.

%%%%%%%%%%%%%%%%%%%%%%%%%%%%%%%%%%%%%%%%%%%%%%%%%%%%%%%%%%%%
\subsubsection{Ignoring Hidden Constraints}

Differentiating

\[
g(x)=0
\]

can reveal required velocity and acceleration conditions.

Satisfying only the position-level constraint may be insufficient.

%%%%%%%%%%%%%%%%%%%%%%%%%%%%%%%%%%%%%%%%%%%%%%%%%%%%%%%%%%%%
\subsubsection{Differentiating Constraints and Discarding the Original Equation}

The derivative condition

\[
\dot g=0
\]

only preserves whatever value of \(g\) is initially present.

The original constraint

\[
g=0
\]

must still be enforced.

%%%%%%%%%%%%%%%%%%%%%%%%%%%%%%%%%%%%%%%%%%%%%%%%%%%%%%%%%%%%
\subsubsection{Assuming the Mass Matrix Is Invertible}

If

\[
M
\]

is singular, one cannot simply write

\[
x'=M^{-1}f.
\]

The singular structure must be analyzed.

%%%%%%%%%%%%%%%%%%%%%%%%%%%%%%%%%%%%%%%%%%%%%%%%%%%%%%%%%%%%
\subsubsection{Using Arbitrary ODE Solvers on Higher-Index DAEs}

Standard explicit ODE methods may violate constraints or fail entirely.

DAE-aware methods are usually necessary.

%%%%%%%%%%%%%%%%%%%%%%%%%%%%%%%%%%%%%%%%%%%%%%%%%%%%%%%%%%%%
\subsubsection{Confusing Algebraic Variables with Constants}

An algebraic variable can change with time.

It is called algebraic because its value is determined through an
algebraic relation rather than an independent differential equation.

%%%%%%%%%%%%%%%%%%%%%%%%%%%%%%%%%%%%%%%%%%%%%%%%%%%%%%%%%%%%
\subsubsection{Ignoring Constraint Drift}

A numerical trajectory can appear visually reasonable while slowly
violating

\[
g(x)=0.
\]

Constraint residuals should be monitored directly.

%%%%%%%%%%%%%%%%%%%%%%%%%%%%%%%%%%%%%%%%%%%%%%%%%%%%%%%%%%%%
\subsubsection{Ignoring Scaling}

Poorly scaled equations can make Newton iterations or linear systems
appear nearly singular.

Physical units and variable magnitudes matter numerically.

%%%%%%%%%%%%%%%%%%%%%%%%%%%%%%%%%%%%%%%%%%%%%%%%%%%%%%%%%%%%
\subsubsection{Confusing DAE Index Definitions}

Differentiation index, tractability index, perturbation index, and other
definitions are related but not identical.

A stated index should be tied to a specific definition.

%%%%%%%%%%%%%%%%%%%%%%%%%%%%%%%%%%%%%%%%%%%%%%%%%%%%%%%%%%%%
\subsubsection{Assuming Higher Index Means No Solution Exists}

Higher index means the constraints are more deeply coupled to the
differential structure.

Solutions can exist, but initialization and numerical treatment become
more delicate.

%%%%%%%%%%%%%%%%%%%%%%%%%%%%%%%%%%%%%%%%%%%%%%%%%%%%%%%%%%%%
\subsubsection{Ignoring Events That Change Constraints}

When contact begins or a circuit switch changes position, the active DAE
may change.

The post-event state must satisfy the new constraints.

%%%%%%%%%%%%%%%%%%%%%%%%%%%%%%%%%%%%%%%%%%%%%%%%%%%%%%%%%%%%
\subsection{Applications}
\label{subsec:dae-applications}
%%%%%%%%%%%%%%%%%%%%%%%%%%%%%%%%%%%%%%%%%%%%%%%%%%%%%%%%%%%%

\begin{applicationbox}

\textbf{Multibody Mechanics.}
Rigid links, joints, pendulums, gears, and contact constraints naturally
lead to differential-algebraic systems.

\medskip

\textbf{Electrical Circuits.}
Kirchhoff laws combine algebraically with capacitor and inductor dynamics.

\medskip

\textbf{Power Systems.}
Generator dynamics interact with algebraic network equations for voltage
and power flow.

\medskip

\textbf{Chemical Engineering.}
Mass and energy balances interact with equilibrium, constitutive, and
phase constraints.

\medskip

\textbf{Robotics.}
Closed kinematic chains and constrained mechanisms generate DAEs with
Lagrange multipliers.

\medskip

\textbf{Control Theory.}
Descriptor systems generalize standard state-space equations to singular
mass matrices.

\medskip

\textbf{Optimization.}
Dynamic optimization and constrained optimal control often produce
differential-algebraic necessary conditions.

\medskip

\textbf{Scientific Simulation.}
Large multiphysics models frequently combine differential laws with
instantaneous algebraic constraints.

\end{applicationbox}

%%%%%%%%%%%%%%%%%%%%%%%%%%%%%%%%%%%%%%%%%%%%%%%%%%%%%%%%%%%%
\subsection{Exercises}
\label{subsec:dae-exercises}
%%%%%%%%%%%%%%%%%%%%%%%%%%%%%%%%%%%%%%%%%%%%%%%%%%%%%%%%%%%%

\begin{exercise}[1]
Define a differential-algebraic equation.
\end{exercise}

\begin{exercise}[1]
Distinguish an ODE from a DAE.
\end{exercise}

\begin{exercise}[1]
Define a differential variable.
\end{exercise}

\begin{exercise}[1]
Define an algebraic variable.
\end{exercise}

\begin{exercise}[1]
Define a consistent initial condition.
\end{exercise}

\begin{exercise}[1]
Define the differentiation index conceptually.
\end{exercise}

\begin{exercise}[1]
Define a mass-matrix DAE.
\end{exercise}

\begin{exercise}[1]
Define a matrix pencil.
\end{exercise}

\begin{exercise}[2]
Identify the differential and algebraic variables in

\[
y'=-y+z,
\]

\[
0=y+z-1.
\]
\end{exercise}

\begin{exercise}[2]
Determine whether

\[
y(0)=2,
\qquad
z(0)=-1
\]

is consistent with

\[
y+z=1.
\]
\end{exercise}

\begin{exercise}[2]
Find the consistent values of \(y(0)\) if

\[
x(0)=\frac12
\]

and

\[
x^2+y^2=1.
\]
\end{exercise}

\begin{exercise}[2]
Differentiate the constraint

\[
x^2+y^2=4.
\]
\end{exercise}

\begin{exercise}[2]
Differentiate the same constraint twice and obtain the acceleration-level
condition.
\end{exercise}

\begin{exercise}[2]
Determine whether the matrix

\[
M
=
\begin{pmatrix}
1&0\\
0&0
\end{pmatrix}
\]

is singular.
\end{exercise}

\begin{exercise}[2]
Reduce

\[
y'=2y+z,
\]

\[
0=y-z
\]

to an ODE.
\end{exercise}

\begin{exercise}[3]
Use the implicit function theorem to explain why nonsingular \(g_z\)
allows local elimination of algebraic variables.
\end{exercise}

\begin{exercise}[3]
Derive the hidden derivative constraint associated with

\[
g(x(t))=0.
\]
\end{exercise}

\begin{exercise}[3]
For the system

\[
x'=v,
\]

\[
v'=\lambda,
\]

\[
x=0,
\]

derive the hidden constraints.
\end{exercise}

\begin{exercise}[3]
Solve the DAE

\[
x'=-x+y,
\]

\[
x+y=1
\]

for consistent initial data.
\end{exercise}

\begin{exercise}[3]
Write the previous DAE in mass-matrix form.
\end{exercise}

\begin{exercise}[3]
Derive the generalized eigenvalue equation for

\[
Ex'=Ax.
\]
\end{exercise}

\begin{exercise}[3]
Explain why singular \(E\) creates algebraic structure.
\end{exercise}

\begin{exercise}[3]
Formulate the Cartesian pendulum as a DAE.
\end{exercise}

\begin{exercise}[3]
Derive the velocity constraint for the Cartesian pendulum.
\end{exercise}

\begin{exercise}[3]
Derive the acceleration constraint for the Cartesian pendulum.
\end{exercise}

\begin{exercise}[3]
Solve algebraically for the pendulum constraint multiplier using the
acceleration-level constraint.
\end{exercise}

\begin{exercise}[3]
Explain how Kirchhoff's current law can create algebraic equations in a
circuit model.
\end{exercise}

\begin{exercise}[3]
Construct a simple capacitor-resistor circuit DAE.
\end{exercise}

\begin{exercise}[3]
Explain how chemical-equilibrium relations can create DAEs.
\end{exercise}

\begin{exercise}[3]
Derive the equilibrium condition

\[
(I-A)x^*=b
\]

for a linear algebraically constrained system where appropriate.
\end{exercise}

\begin{exercise}[3]
Derive the implicit Euler equations for a semi-explicit index-1 DAE.
\end{exercise}

\begin{exercise}[3]
Apply one implicit Euler step to

\[
y'=-y+z,
\]

\[
y+z=1.
\]
\end{exercise}

\begin{exercise}[3]
Explain why Newton's method is needed in nonlinear implicit DAE
integration.
\end{exercise}

\begin{exercise}[3]
Explain constraint drift in a numerical pendulum simulation.
\end{exercise}

\begin{exercise}[3]
Describe a projection step that restores

\[
x^2+y^2=L^2.
\]
\end{exercise}

\begin{exercise}[3]
Explain why differentiating a constraint can enlarge the solution set.
\end{exercise}

\begin{exercise}[3]
Show how the singular perturbation system

\[
y'=f(y,z),
\]

\[
\varepsilon z'=g(y,z)
\]

formally becomes a DAE as

\[
\varepsilon\to0.
\]
\end{exercise}

\begin{exercise}[4]
Study the Weierstrass canonical form for regular matrix pencils.
\end{exercise}

\begin{exercise}[4]
Relate nilpotent blocks in the Weierstrass form to DAE index.
\end{exercise}

\begin{exercise}[4]
Study Kronecker canonical form for singular matrix pencils.
\end{exercise}

\begin{exercise}[4]
Investigate tractability and perturbation indices.
\end{exercise}

\begin{exercise}[4]
Study existence and uniqueness theory for nonlinear DAEs.
\end{exercise}

\begin{exercise}[4]
Develop consistent-initialization algorithms for nonlinear DAEs.
\end{exercise}

\begin{exercise}[4]
Study index reduction by constraint differentiation.
\end{exercise}

\begin{exercise}[4]
Investigate structural index methods such as graph-based DAE analysis.
\end{exercise}

\begin{exercise}[4]
Study constrained mechanical systems using Lagrange multipliers.
\end{exercise}

\begin{exercise}[4]
Analyze the index of the Cartesian pendulum DAE.
\end{exercise}

\begin{exercise}[4]
Study Baumgarte constraint stabilization.
\end{exercise}

\begin{exercise}[4]
Investigate projection methods for constrained mechanics.
\end{exercise}

\begin{exercise}[4]
Study backward differentiation formulas for DAEs.
\end{exercise}

\begin{exercise}[4]
Analyze convergence of implicit Euler for index-1 DAEs.
\end{exercise}

\begin{exercise}[4]
Study numerical difficulties in higher-index DAEs.
\end{exercise}

\begin{exercise}[4]
Investigate descriptor-system controllability and observability.
\end{exercise}

\begin{exercise}[4]
Study generalized eigenvalue stability for descriptor systems.
\end{exercise}

\begin{exercise}[4]
Investigate DAEs arising from modified nodal analysis of electrical
circuits.
\end{exercise}

\begin{exercise}[4]
Study quasi-steady-state reduction and singular perturbations.
\end{exercise}

\begin{exercise}[4]
Investigate hybrid DAEs with switching constraints.
\end{exercise}

\begin{exercise}[4]
Study dynamic optimization with DAE constraints.
\end{exercise}

%%%%%%%%%%%%%%%%%%%%%%%%%%%%%%%%%%%%%%%%%%%%%%%%%%%%%%%%%%%%
\subsection{Conceptual Summary}
\label{subsec:dae-summary}
%%%%%%%%%%%%%%%%%%%%%%%%%%%%%%%%%%%%%%%%%%%%%%%%%%%%%%%%%%%%

A differential-algebraic equation combines differential evolution with
algebraic constraints.

A general implicit form is

\[
\boxed{
F(t,x,x')=0.
}
\]

A semi-explicit DAE has form

\[
\boxed{
\begin{aligned}
y'&=f(t,y,z),\\
0&=g(t,y,z).
\end{aligned}
}
\]

Consistent initial conditions satisfy

\[
\boxed{
g(t_0,y_0,z_0)=0
}
\]

together with any necessary hidden constraints.

A mass-matrix system is

\[
\boxed{
M x'=f(t,x).
}
\]

If \(M\) is singular, algebraic constraints may be present.

A linear DAE is

\[
\boxed{
E x'=Ax+f(t),
}
\]

with matrix pencil

\[
\boxed{
\lambda E-A.
}
\]

Higher-index systems may require differentiation of algebraic constraints
before explicit derivatives can be determined.

For numerical integration, implicit methods solve differential and
algebraic equations simultaneously.

Thus DAEs combine

\[
\boxed{
\text{differential equations}
+
\text{algebraic constraints}
+
\text{linear algebra}
+
\text{numerical analysis}
+
\text{constrained dynamics}.
}
\]

%%%%%%%%%%%%%%%%%%%%%%%%%%%%%%%%%%%%%%%%%%%%%%%%%%%%%%%%%%%%
\subsection{Section Connections}
\label{subsec:dae-connections}
%%%%%%%%%%%%%%%%%%%%%%%%%%%%%%%%%%%%%%%%%%%%%%%%%%%%%%%%%%%%

\chapterconnection{
Differential-Algebraic Equations complete Chapter~9 by extending
differential equations to systems subject to instantaneous algebraic
constraints. They connect ordinary differential equations with
constrained mechanics, electrical circuits, optimization, singular
perturbation theory, and numerical simulation. Their mass-matrix and
generalized-eigenvalue formulations link directly with linear algebra,
while higher-index constraints reveal why specialized numerical methods
are required. Chapter 10 develops Optimization and Operations Research,
where constrained optimization, dynamic programming, stochastic
optimization, networks, scheduling, decision theory, game theory, and
optimal control extend many of the structural ideas encountered here.
}

%%%%%%%%%%%%%%%%%%%%%%%%%%%%%%%%%%%%%%%%%%%%%%%%%%%%%%%%%%%%
% SECTION SOURCE TRACKING
%%%%%%%%%%%%%%%%%%%%%%%%%%%%%%%%%%%%%%%%%%%%%%%%%%%%%%%%%%%%

%------------------------------------------------------------
% 9.8 Differential-Algebraic Equations
%------------------------------------------------------------
%
% Source basis:
%   - Standard differential-algebraic equation theory.
%   - Standard constrained mechanical and descriptor-system models.
%   - Standard numerical treatment of index-1 and higher-index DAEs.
%
% Used for:
%   - DAE definitions;
%   - implicit and semi-explicit formulations;
%   - differential and algebraic variables;
%   - algebraic constraints;
%   - constraint manifolds;
%   - consistent initial conditions;
%   - hidden constraints;
%   - tangent-space conditions;
%   - differentiation index;
%   - index-1 reduction;
%   - higher-index systems;
%   - mass-matrix form;
%   - linear constant-coefficient DAEs;
%   - matrix pencils;
%   - generalized eigenvalues;
%   - constrained mechanical systems;
%   - Lagrange multipliers;
%   - Cartesian pendulum DAEs;
%   - circuit DAEs;
%   - chemical and process-system DAEs;
%   - conservation constraints;
%   - descriptor control systems;
%   - consistent initialization;
%   - index reduction;
%   - constraint drift;
%   - projection methods;
%   - Baumgarte stabilization preview;
%   - implicit Euler;
%   - backward differentiation formulas;
%   - Newton iteration;
%   - residual formulations;
%   - stiffness and scaling;
%   - events and hybrid DAEs;
%   - singular perturbation connections;
%   - applications and exercises.
%
% Notes:
%   - This completes Chapter 9 — Differential Equations and Dynamical
%     Systems.
%   - Chapter 10 begins Optimization and Operations Research.
%   - Numerical DAE methods are developed more deeply in Chapter 11.
%
%%%%%%%%%%%%%%%%%%%%%%%%%%%%%%%%%%%%%%%%%%%%%%%%%%%%%%%%%%%%